%%%%%%%%%%%%%%%%%%%%%%%%%%%%%%%%%%%%%%%%%%%%%%%%%%%%%%%%%%%%%%%
%Documento con el glosario de términos. Si el glosario de términos es muy largo y acaba en pagina par comentar el comando \blankpage y decomentar el comando \newpage
%%%%%%%%%%%%%%%%%%%%%%%%%%%%%%%%%%%%%%%%%%%%%%%%%%%%%%%%%%%%%%%
\begin{center}
    {\fontsize{20}{80}\selectfont \textbf{Glosario de Términos}\\}
\end{center}
\bigskip

\begin{description}
    \item[Bloque] Un bloque en Simulink es una entidad gráfica que representa una operación, función o elemento específico dentro de un modelo, utilizado para construir y simular sistemas dinámicos.
    \item[Controlador] Programa que permite a una computadora u otro dispositivo electrónico manejar los componentes que tiene instalados.
    \item[Controlador de Potencia] Dispositivo que regula y gestiona el flujo de energía eléctrica a un sistema o componente, asegurando un funcionamiento eficiente y seguro.
    \item[Datos Estándar] Los Datos Estándar son aquellos que ocupan 11 bits en el protocolo CAN y se utilizan principalmente para la comunicación entre dispositivos en sistemas automotrices.
    \item[Datos Extendidos] Los Datos Extendidos ocupan 29 bits y se emplean para transmitir datos en redes más complejas y grandes, como en aplicaciones industriales y de automatización.
    \item[Firmware] Software integrado en dispositivos electrónicos que controla sus funciones básicas y proporciona una interfaz entre el hardware y el software.
    \item[Hardware] Conjunto de aparatos de una computadora.
    \item[Microcontrolador] Pequeño ordenador integrado en un solo chip que contiene un procesador, memoria y periféricos, utilizado para controlar dispositivos electrónicos.
    \item[Middleware] Software que actúa como intermediario facilitando la comunicación y gestión de datos entre diferentes aplicaciones o sistemas.
    \item[Modelo] Un modelo en Simulink es una representación visual de un sistema dinámico o proceso, creado mediante bloques funcionales interconectados que describen su funcionamiento y relaciones.
    \item[Motor BLDC] Motor eléctrico de corriente continua sin escobillas.
    \item[Paquetes Informáticos] Conjunto de programas diseñados para realizar tareas específicas en un ordenador.
    \item[Periférico] Aparato auxiliar e independiente conectado a la unidad central de una computadora u otro dispositivo electrónico.
    \item[Servomotores] Sistema electromecánico que amplifica la potencia reguladora.
    \item[Shield] Módulo de expansión que se conecta a una placa para añadirle funcionalidades adicionales, como conectividad, sensores o controladores.
    \item[Software] Conjunto de programas, instrucciones y reglas informáticas para ejecutar ciertas tareas en una computadora.
\end{description}

%\blankpage
%\newpage